下一个有趣的问题当然是 $\mathcal{A}$ 的指标是否有限, 更具体地说, 何时为零. 回答这一问题的标准方法之一是利用 Fredholm 指标的同伦不变性. 例如, 考虑二阶椭圆算子
\begin{equation}\label{eq:elliptic-dirichlet-index-homotopy-operator}
L(\bm{x},D)u=a_{ij}(\bm{x})\frac{\partial^2u}{\partial x_i\partial x_j}+b_i(\bm{x})\frac{\partial u}{\partial x_i}+c(\bm{x})u
\end{equation}
及 Dirichlet 边界条件 $B(\bm{x},D)u=u$. 假设矩阵 $a_{ij}$ 对称且正定. 于是可以考虑算子的单参数族
\begin{equation}\label{eq:elliptic-dirichlet-laplace-homotopy}
L_t=(1-t)\Delta+tL,\quad B_t=B.
\end{equation}


若 $\Omega$ 和系数满足相应的光滑性假设, 则对每个 $t\in[0,1]$, Theorem~\ref{thm:agmon-douglis-nirenberg-elliptic-estimate} 的假设均适用; 因而 $(L,B)$ 的 Fredholm 指标与 Laplace 方程的相同. 在 Section~\ref{sec:dirichlet-existence-and-uniqueness} 中已经证明, 对每个 $f\in H^{-1}(\Omega)$, 问题 $\Delta u=f$ 在边界条件 $u=0$ 下有唯一解 $u\in H^1(\Omega)$. 利用逆迹定理, 容易得到对每个 $f\in H^{-1}(\Omega)$, $g\in H^{1/2}(\Omega)$, 问题 $\Delta u=f$, $u|_{\partial\Omega}=g$ 有唯一解 $u\in H^1(\Omega)$. 现在希望知道, 若 $f\in L^2(\Omega)$ 且 $g\in H^{3/2}(\Omega)$, 是否就有 $u\in H^2(\Omega)$. 这与 Theorem~\ref{thm:agmon-douglis-nirenberg-elliptic-estimate} 的第一项断言十分相近, 但实际上并不能由 Theorem~\ref{thm:agmon-douglis-nirenberg-elliptic-estimate} 推出. 原因是, 对 Laplace 方程的 Dirichlet 问题, 应选取 $s_1=0$, $t_1=2$ 和 $r_1=-2$, 从而 $M_1=0$ 且 $X_{M_1}=H^2(\Omega)$. 因此, 该定理断言数据适当时 $H^2$ 解具有更高正则性, 而不是断言 $H^1$ 解具有 $H^2$ 正则性. 尽管如此, 可以沿着与 Theorem~\ref{thm:agmon-douglis-nirenberg-elliptic-estimate} 非常相似的思路证明弱解的正则性, 而 Agmon, Douglis 和 Nirenberg 实际上也对标量椭圆方程陈述了这样的结果. 具有 Dirichlet 条件的二阶方程参见 Sections~\ref{sec:elliptic-interior-regularity} 和 \ref{sec:elliptic-boundary-regularity}.

现在自然要问, 基于 Lax--Milgram 引理的 Section~\ref{sec:dirichlet-existence-and-uniqueness} 中的方法可以推广到哪一类问题. 这将引向 Subsection~\ref{subsec:agmon-condition-coercive-problems} 将讨论的 Agmon 条件. Lax--Milgram 引理将给出一个 ``弱\textquotesingle\textquotesingle\ 解的存在性, 而在应用 Theorem~\ref{thm:agmon-douglis-nirenberg-elliptic-estimate} 之前, 仍然必须处理弱解的正则性.

另一个有趣的问题是刻画 $\mathcal{A}$ 的值域的正交补; 即 $\bm{f}$ 和 $\bm{g}$ 必须满足什么条件, 问题 $\bm{L}\bm{u}=\bm{f}$ 才能在边界条件 $\bm{B}\bm{u}=\bm{g}$ 下可解? 通常, 应用逆迹定理就能找到满足 $\bm{B}\bm{u}=\bm{g}$ 的 $\bm{u}$ (见下一小节, Subsection~\ref{subsec:adjoint-boundary-value-problem}); 因而归结为 $\bm{g}=0$ 的情形. 于是剩下的问题是刻画这样的 $\bm{v}$: 对每个满足 $\bm{B}\bm{u}=0$ 的 $\bm{u}$, 都有 $(\bm{v},\bm{L}\bm{u})=0$. 形式地分部积分可以得到关于 $\bm{v}$ 的椭圆边值问题, 称为伴随边值问题. 下一小节将研究标量椭圆方程的伴随边值问题. 当然, 先验而言, $\bm{v}$ 只在 ``弱\textquotesingle\textquotesingle\ 或 ``广义\textquotesingle\textquotesingle\ 意义下满足伴随边值问题. 因此弱解的正则性再次成为重要问题. 特别地, 为证明算子 $\mathcal{A}$ 为 Fredholm 算子, 必须证明伴随算子的零空间有限维. 当然, 必须对伴随问题的弱解证明这一点, 而不能仅对强解证明. 事实上可以证明这一点. 若系数足够光滑, 则伴随问题的弱解实际上是光滑的.
